\documentclass[../../../include/open-logic-section]{subfiles}

\begin{document}

% Part: sets-functions-relations
% Chapter: sets
% Section: russells-paradox

\olfileid{sfr}{set}{rus}
\olsection{Russell's Paradox}

外延性原理为记号 $\Setabs{x}{\phi(x)}$ 提供了依据，该记号表示满足~$\phi(x)$ 的那些~$x$ 所构成的\emph{那个}集合。然而，外延性原理\emph{真正}认可的只是下面这个论断：\emph{如果}存在一个集合，其成员恰好是且仅是所有具有性质~$\phi$ 的对象，\emph{那么}这样的集合只有一个。换句话说：固定某个~$\phi$ 后，集合 $\Setabs{x}{\phi(x)}$ 是唯一的——\emph{倘若它存在的话}。

但这个条件句非常重要！关键之处在于，并非每个性质都能用于\emph{概括}。也就是说，有些性质\emph{并不}定义集合。如果所有性质都能定义集合，那么我们就会陷入直接的矛盾。其中最著名的例子就是 Russell（罗素）悖论。

集合可以是其他集合的元素——例如，一个集合~$A$ 的幂集就是由集合构成的。因此，询问或研究一个集合是否是另一个集合的元素是有意义的。一个集合能是自己的元素吗？从集合的观念来看，似乎没有什么能排除这种可能。例如，如果\emph{所有}集合形成一个对象的汇集，人们可能会认为它们可以被归拢为一个单一的集合——即所有集合的集合。而这个集合，作为一个集合，将会是所有集合的集合的一个元素。

当我们考虑“不属于自身”这一性质，即\emph{非自属的}性质时，罗素 悖论就产生了。如果我们假设存在一个由所有不属于自身的集合构成的集合，会怎样呢？
\[
R = \Setabs{x}{x \notin x}
\]
存在吗？结果表明，我们可以证明它并不存在。

\begin{thm}[Russell 悖论]\ollabel{thm:russells-paradox}
	不存在集合 $R = \Setabs{x}{x \notin x}$。
\end{thm}

\begin{proof}
如果 $R = \Setabs{x}{x \notin x}$ 存在，那么
$R \in R$ 当且仅当 $R \notin R$，这是一个矛盾。
\end{proof}

\begin{tagblock}{novice}
\begin{explain}
让我们更详细地梳理一下这个证明。如果 $R$ 存在，那么考察 $R \in R$ 是否成立是有意义的。假设确实 $R \in R$。现在，$R$ 被定义为所有不属于自身的集合的集合。所以，如果 $R \in R$，那么 $R$ 本身就不具有 $R$ 的定义性质。但只有具有该性质的集合才在~$R$ 中，因此，$R$ 不可能是~$R$ 的元素，即 $R \notin R$。但 $R$ 不可能同时是且不是~$R$ 的元素，所以我们得到了一个矛盾。

既然假设 $R \in R$ 会导致矛盾，我们就有 $R \notin R$。但这也导致矛盾！因为如果 $R
\notin R$，那么 $R$ 本身确实具有 $R$ 的定义性质，因此 $R$ 将像所有其他非自属集合一样，是 $R$ 的一个元素。同样，它不可能同时不是且是~$R$ 的元素。
\end{explain}
\end{tagblock}

\begin{digress}
我们如何建立一个能避免陷入 罗素 悖论的集合论，即避免做出 $R = \Setabs{x}{x \notin x}$ 存在这一\emph{不一致的}断言呢？我们需要设立一些公理，为我们提供陈述集合何时存在（以及何时不存在）的非常精确的条件。
	
本章概述的集合论没有做到这一点。它是\emph{真正朴素的}。它只告诉你集合服从外延性，并且如果你有一些集合，你可以构造它们的并集、交集等。我们有可能以比这更严格的方式发展集合论。 \oliflabeldef{cumul:::part}{这种严格性将留到第 \olref[cumul][][]{part} 部分。目前，我们将在朴素的意义下推进，并小心地尝试避开矛盾。}{}
\end{digress}

\end{document}